\section{Problem formulation}
\label{sec:problem}

\subsection{Safety and utility objectives}

Let $Q$ be a workload of authenticated task contexts and $G(q)\subseteq\M$ the policy ground-truth items entitled and relevant to task $q$. Let $O(q)$ be the items returned by the store. Authorization violations and misses are
\begin{align}
V(Q)&=\sum_{q\in Q}\sum_{m\in O(q)}\mathbf{1}[\neg\Auth(m,q)],\\
L(Q)&=\sum_{q\in Q}|G(q)\setminus O(q)|.
\end{align}
The safety constraint is $V(Q)=0$ for service-mediated reads under the model. Utility seeks to reduce $L(Q)$, answer error, latency, energy, and storage subject to safety and device budgets.

Because relevance and authorization are different predicates, maximizing conventional recall over all historically relevant items is incorrect. We report \emph{entitled recall}
\begin{equation}
\mathsf{ER@}k=\frac{1}{|Q^{+}|}\sum_{q\in Q^{+}}\frac{|O_k(q)\cap G(q)|}{|G(q)|},
\label{eq:er}
\end{equation}
where $Q^{+}=\{q\in Q: G(q)\ne\emptyset\}$ is the set of queries with nonempty entitled ground truth, alongside out-of-scope retrieval rate. The convention is frozen: queries with $G(q)=\emptyset$ contribute to no denominator and are excluded from the reported mean, a choice that depends only on the method-independent ground truth and not on any method's output, so $\mathsf{ER@}k$ stays comparable across methods evaluated on the identical query set. Reported results also state $|Q^{+}|$ and the number of excluded queries $|Q\setminus Q^{+}|$. A method cannot trade a nonzero violation rate for better entitled recall and still satisfy the principal security property.

\subsection{Derived-item safety}

Let $\DAG=(\I,E)$ be the derivation DAG, where $(j,i)\in E$ means item $i$ cites item $j$ as a source. For parent set $\mathrm{pa}(i)$, a derived item is policy-valid when
\begin{equation}
\ell_i\preceq\Meet(\mathrm{pa}(i))
\quad\wedge\quad
\mathsf{Attributed}(x_i,\{x_j:j\in\mathrm{pa}(i)\}).
\label{eq:derived-validity}
\end{equation}
The first conjunct is deterministic. The second is necessarily approximate for natural-language derivations. MemScope treats a false acceptance by the attribution checker as an integrity risk, not as evidence that a broader label is safe. A deployment can require extractive summaries or user confirmation for high-sensitivity classes.

\subsection{Storage-budget objective}

Let $w_i$ be item size and $y_i\in\{0,1\}$ indicate retention after compaction. For unexpired, policy-retainable items, define estimated future utility $u_i$, safety value $r_i$, sensitivity cost $d_i$, staleness $g_i$, and recomputation cost $c_i$. The planner solves or approximates
\begin{align}
\max_{\bm y}\quad &\sum_i y_i(\alpha u_i+\beta r_i+\gamma c_i-\delta d_i-\eta g_i)\\
\text{s.t.}\quad &\sum_i y_iw_i\le B,\quad y_i=0\text{ if }\tau\ge e_i,
\label{eq:eviction}
\end{align}
with policy-pinned items handled as hard constraints only when their total size fits $B$. If hard constraints exceed the budget, the service raises an explicit resource fault rather than silently discarding safety-critical records. The objective affects retention, never authorization.

\subsection{Design requirements}

The formulation yields seven requirements: (R1) authenticated principal and purpose binding; (R2) per-document enforcement before ranking; (R3) conservative derivation; (R4) retrieval-time expiry; (R5) atomic policy/version checks; (R6) auditable lifecycle transitions; and (R7) bounded mobile overhead. The next sections instantiate these requirements.
